%% Created by Maple 2024.2, Linux
%% Source Worksheet: Laborator_1.mw
%% Generated: Wed Apr 15 22:35:49 EEST 2026
\documentclass{article}
\usepackage{amssymb}
\usepackage{graphicx}
\usepackage{hyperref}
\usepackage{listings}
\usepackage{mathtools}
\usepackage{maple}
\usepackage[utf8]{inputenc}
\usepackage[svgnames]{xcolor}
\usepackage{amsmath}
\usepackage{breqn}
\usepackage{textcomp}
\begin{document}
\lstset{basicstyle=\ttfamily,breaklines=true,columns=flexible}
\pagestyle{empty}
\DefineParaStyle{Maple Bullet Item}
\DefineParaStyle{Maple Heading 1}
\DefineParaStyle{Maple Warning}
\DefineParaStyle{Maple Heading 4}
\DefineParaStyle{Maple Heading 2}
\DefineParaStyle{Maple Heading 3}
\DefineParaStyle{Maple Dash Item}
\DefineParaStyle{Maple Error}
\DefineParaStyle{Maple Title}
\DefineParaStyle{Maple Ordered List 1}
\DefineParaStyle{Maple Text Output}
\DefineParaStyle{Maple Ordered List 2}
\DefineParaStyle{Maple Ordered List 3}
\DefineParaStyle{Maple Normal}
\DefineParaStyle{Maple Ordered List 4}
\DefineParaStyle{Maple Ordered List 5}
\DefineCharStyle{Maple 2D Output}
\DefineCharStyle{Maple 2D Input}
\DefineCharStyle{Maple Maple Input}
\DefineCharStyle{Maple 2D Math}
\DefineCharStyle{Maple Hyperlink}
\mapleinput
{$ \displaystyle \texttt{>\,} 12+4-5 $}

% \mapleresult
\begin{dmath}\label{(1)}
11
\end{dmath}
\mapleinput
{$ \displaystyle \texttt{>\,} 2^{10} $}

% \mapleresult
\begin{dmath}\label{(2)}
1024
\end{dmath}
\mapleinput
{$ \displaystyle \texttt{>\,} \sin ( 0.1) $}

% \mapleresult
\begin{dmath}\label{(3)}
 0.09983341665
\end{dmath}
\mapleinput
{$ \displaystyle \texttt{>\,} (a +b)(a -b) $}

% \mapleresult
\begin{dmath}\label{(4)}
a \! \left(a -b \right)+b \! \left(a -b \right)
\end{dmath}
\mapleinput
{$ \displaystyle \texttt{>\,} f \coloneqq x \rightarrow 3\,x^{3}+2\,x^{2}-5\colon  $}

\mapleinput
{$ \displaystyle \texttt{>\,} \mathit{diff} (f (x),x) $}

% \mapleresult
\begin{dmath}\label{(5)}
9 x^{2}+4 x 
\end{dmath}
\mapleinput
{$ \displaystyle \texttt{>\,} f \coloneqq x \rightarrow \mathrm{sqrt}(1+x^{4})\colon  $}

\mapleinput
{$ \displaystyle \texttt{>\,} \mathit{diff} (f (x),x) $}

% \mapleresult
\begin{dmath}\label{(6)}
\frac{2 x^{3}}{\sqrt{x^{4}+1}}
\end{dmath}
\mapleinput
{$ \displaystyle \texttt{>\,} f \coloneqq x \rightarrow \exp (x)\cdot \sin (x)\cdot \cos (x)\colon  $}

\mapleinput
{$ \displaystyle \texttt{>\,} \mathit{diff} (f (x),x) $}

% \mapleresult
\begin{dmath}\label{(7)}
{\mathrm e}^{x} \sin \! \left(x \right) \cos \! \left(x \right)+{\mathrm e}^{x} \cos \! \left(x \right)^{2}-{\mathrm e}^{x} \sin \! \left(x \right)^{2}
\end{dmath}
\mapleinput
{$ \displaystyle \texttt{>\,} \mathit{simplify} ((a +b)(a -b)) $}

% \mapleresult
\begin{dmath}\label{(8)}
a \! \left(a -b \right)+b \! \left(a -b \right)
\end{dmath}
\mapleinput
{$ \displaystyle \texttt{>\,} \mathit{expand} ((a +b)\cdot (a -b)) $}

% \mapleresult
\begin{dmath}\label{(9)}
a^{2}-b^{2}
\end{dmath}
\mapleinput
{$ \displaystyle \texttt{>\,} \mathit{simplify} ((a -b)\cdot (a +b)) $}

% \mapleresult
\begin{dmath}\label{(10)}
a^{2}-b^{2}
\end{dmath}
\mapleinput
{$ \displaystyle \texttt{>\,} \mathit{int} (3\,x^{3}+2\,x^{2}-5,\,x =0\,..\,1) $}

% \mapleresult
\begin{dmath}\label{(11)}
-\frac{43}{12}
\end{dmath}
\mapleinput
{$ \displaystyle \texttt{>\,} \mathit{int} (\frac{1}{x^{2}},\,x =0..\mathrm{infinity}) $}

% \mapleresult
\begin{dmath}\label{(12)}
\infty 
\end{dmath}
\mapleinput
{$ \displaystyle \texttt{>\,} \mathit{int} (\exp (-x^{2}),\,x =-\mathrm{infinity}..\,\mathrm{infinity}); $}

% \mapleresult
\begin{dmath}\label{(13)}
\sqrt{\pi}
\end{dmath}
\mapleinput
{$ \displaystyle \texttt{>\,} \mathit{limit} (\frac{\sin (x)}{x},x =0) $}

% \mapleresult
\begin{dmath}\label{(14)}
1
\end{dmath}
\mapleinput
{$ \displaystyle \texttt{>\,} \mathit{limit} (\frac{(x^{3}+3\,x^{2}-5)}{2\,x^{3}-7\,x},x =\mathrm{infinity}) $}

% \mapleresult
\begin{dmath}\label{(15)}
\frac{1}{2}
\end{dmath}
\mapleinput
{$ \displaystyle \texttt{>\,} \mathit{limit} (\frac{(\cos (x)+1)}{x -\mathrm{Pi}},x =\mathrm{Pi}) $}

% \mapleresult
\begin{dmath}\label{(16)}
0
\end{dmath}
\mapleinput
{$ \displaystyle \texttt{>\,} f \coloneqq x \rightarrow \exp (-x)-1\colon  $}

\mapleinput
{$ \displaystyle \texttt{>\,} \mathit{with} (\mathit{plots})\colon  $}

\mapleinput
{$ \displaystyle \texttt{>\,} \mathit{plot} (f (x),x =-2..2) $}

% \mapleresult
\mapleplot{Laborator_2plot2d1.png}
\mapleinput
{$ \displaystyle \texttt{>\,} f \coloneqq (x ,r)\rightarrow \frac{(200\cdot \exp (r \cdot x))}{2(\exp (r \cdot x)-1)+100}\colon  $}

\mapleinput
{$ \displaystyle \texttt{>\,} \mathit{with} (\mathit{plots})\colon  $}

\mapleinput
{$ \displaystyle \texttt{>\,} \mathit{plot} ([f (x , 0.5),f (x ,- 0.5)],x =0..50),\mathit{color} =\,[\mathit{red} ,\mathit{blue}] $}

% \mapleresult
\mapleinput
{$ \displaystyle \texttt{>\,} f \coloneqq x \rightarrow x \cdot \sin (\frac{1}{x})\colon  $}

\mapleinput
{$ \displaystyle \texttt{>\,} \mathit{with} (\mathit{plots})\colon  $}

\mapleinput
{$ \displaystyle \texttt{>\,} \mathit{plot} (f (x),x =-3..3) $}

% \mapleresult
\mapleplot{Laborator_2plot2d2.png}
\mapleinput
{$ \displaystyle \texttt{>\,} \mathit{with} (\mathit{plots})\colon  $}

\mapleinput
{$ \displaystyle \texttt{>\,} \mathit{plot} ([(1-\cos (t)\cdot \cos (t)),\,(1-\cos (t))\cdot \sin (t),\,t =0..2\,\mathrm{Pi}]) $}

% \mapleresult
\mapleplot{Laborator_2plot2d3.png}
\mapleinput
{$ \displaystyle \texttt{>\,} \mathit{with} (\mathit{plots})\colon  $}

\mapleinput
{$ \displaystyle \texttt{>\,} \mathit{plot} ([\sin (3\cdot t)\cos (t),\,\sin (3\cdot t)\sin (t),\,t =0..2\cdot \mathrm{Pi}]) $}

% \mapleresult
\mapleplot{Laborator_2plot2d4.png}
\mapleinput
{$ \displaystyle \texttt{>\,} \mathit{with} (\mathit{plots})\colon  $}

\mapleinput
{$ \displaystyle \texttt{>\,} \mathit{plot} ([t -\sin (t),1-\cos (t),t =0..6\cdot \mathrm{Pi}\,]) $}

% \mapleresult
\mapleplot{Laborator_2plot2d5.png}
\mapleinput
{$ \displaystyle \texttt{>\,} f \coloneqq (t ,s)\rightarrow 1-(\frac{s \cdot \cos (4\,t)\cdot \cos (t)}{\mathrm{sqrt}(1-s^{2}\cdot \cos^{2}(4\,t)\cdot \sin^{2}(t))})\colon  $}

\mapleinput
{$ \displaystyle \texttt{>\,} x \coloneqq (t ,s)\rightarrow f (t -\frac{\mathrm{Pi}}{2},s)\colon  $}

\mapleinput
{$ \displaystyle \texttt{>\,} y \coloneqq (t ,s)\rightarrow f (t ,s)\colon  $}

\mapleinput
{$ \displaystyle \texttt{>\,} \textit{list\_xy} \coloneqq [x (t ,\frac{s}{10}),y (t ,\frac{s}{10}),t =0..2\cdot \mathrm{Pi}]\$ s =1..10\colon  $}

\mapleinput
{$ \displaystyle \texttt{>\,} \mathit{plot} ([\textit{list\_xy}]) $}

% \mapleresult
\mapleplot{Laborator_2plot2d6.png}
\mapleinput
{$ \displaystyle \texttt{>\,} \mathit{with} (\mathit{plots})\colon  $}

\mapleinput
{$ \displaystyle \texttt{>\,} \mathit{implicitplot} (x^{2}+y^{2}-2\,x -4\,y +4=0) $}

% \mapleresult
\mapleplot{Laborator_2plot2d7.png}
\mapleinput
{$ \displaystyle \texttt{>\,} \mathit{with} (\mathit{plots})\colon  $}

\mapleinput
{$ \displaystyle \texttt{>\,} \mathit{implicitplot} ((x -y)^{2}-2\,y^{2}=1) $}

% \mapleresult
\mapleplot{Laborator_2plot2d8.png}
\mapleinput
{$ \displaystyle \texttt{>\,} z \coloneqq (x ,y)\rightarrow 4\,x^{2}\cdot \exp (y)-2\,x^{4}-\exp (4\,y)\colon  $}

\mapleinput
{$ \displaystyle \texttt{>\,} \mathit{with} (\mathit{plots})\colon  $}

\mapleinput
{$ \displaystyle \texttt{>\,} \mathit{plot3d} (z (x ,y),x =-3..3,\,y =-1..1,\mathit{axes} =\mathit{boxed}) $}

% \mapleresult
\mapleplot{Laborator_2plot3d9.png}
\mapleinput
{$ \displaystyle \texttt{>\,} z \coloneqq (x ,y)\rightarrow 4\,x^{2}-y^{2}\colon  $}

\mapleinput
{$ \displaystyle \texttt{>\,} \mathit{with} (\mathit{plots})\colon  $}

\mapleinput
{$ \displaystyle \texttt{>\,} \mathit{plot3d} (z (x ,y),\,x =-100..100,\,y =-100..100,\,\mathit{axes} =\,\mathit{boxed}) $}

% \mapleresult
\mapleplot{Laborator_2plot3d10.png}
\mapleinput
{$ \displaystyle \texttt{>\,} \mathit{with} (\mathit{linalg})\colon  $}

\mapleinput
{$ \displaystyle \texttt{>\,} A \coloneqq \mathit{matrix} ([[1,2,-1],\,[0,1,0],\,[3,-1,2]])\colon  $}

\mapleinput
{$ \displaystyle \texttt{>\,} B \coloneqq \mathit{matrix} ([[1,2,3],\,[1,1,2],[2,1,1]])\colon  $}

\mapleinput
{$ \displaystyle \texttt{>\,} C \coloneqq \mathit{matrix} ([[2,1,1],[0,1,-1],[4,2,2]])\colon  $}

\mapleinput
{$ \displaystyle \texttt{>\,} \mathit{evalm} (2\,\cdot A -B \:\&\!\!*C); $}

% \mapleresult
\mapleinput
{$ \displaystyle \texttt{>\,} \mathit{evalm} (A^{\mathit{(-1)}});
\\
  $}

% \mapleresult
\mapleinput
{$ \displaystyle \texttt{>\,} \mathit{eigenvects} (C); $}

% \mapleresult
\mapleinput
{$ \displaystyle \texttt{>\,} \mathit{eigenvals} (C) $}

% \mapleresult
\begin{dmath}\label{(20)}
0,3,2
\end{dmath}
\mapleinput
{$ \displaystyle \texttt{>\,} \, $}

\end{document}
